%%%%%%%%%%%%%%%%%%%%%%%%%%%%%%%%%%%%%%%%%%%%%%%%%
\section{Summary}
\label{sec:summary}

We demonstrated that the behavior of UV divergences of quasi-PDFs is very different from that of PDFs.  While the renormalization of quasi-PDFs is a simple multiplicative factor, the renormalization of PDFs is of a convolution form, due to the fact that the UV divergences of PDFs are not completely local, and consequently, PDFs of different flavors mix with each other under the renormalization.

We show that the locality in space-time of the perturbative UV divergences of the coordinate-space quasi-PDFs makes it possible to have the coordinate-space quasi-PDFs multiplicatively renormalizable, as shown in Eq.~(\ref{eq:renor}), to all orders in QCD perturbation theory.  With the all-order arguments for factorizing all leading power CO divergences of quasi-PDFs into PDFs, given in Ref.~\cite{Ma:2014jla}, we conclude that the renormalized quasi-PDFs could be good candidates for extracting PDFs from LQCD calculations.

For the LQCD calculation of quasi-PDFs, however, the introduction of the overall factor, $e^{-C_i |\xi_z|}$, with perturbatively calculated $C_i$, is not sufficient to remove all power divergences, since we cannot calculate the $C_j$ to all orders.  That is, we need to renormalize the power UV divergence of quasi-PDFs, nonperturbatively.  As shown in Ref.~\cite{Ishikawa:2016znu}, for example, we could introduce an overall non-perturbative factor to remove this kind of power divergences to all orders, for which we will not get into the details here.  In principle, the choice to renormalize the power divergences of quasi-PDFs nonperturbatively is not unique.  But, so long as the renormalization is multiplicative, the renormalization procedure does not alter the CO properties of quasi-PDFs (so that the quasi-PDFs could be factorized into PDFs), the difference in renormalization procedures/choices corresponds to different renormalization schemes, which should lead to different matching coefficients between quasi-PDFs and PDFs.

In addition, we show that the coordinate-space quasi-PDFs are well behaved for all values of $\xi_z$, except when $\xi_z=0$.  That is,
if we want to Fourier transform the coordinate-space quasi-PDFs to derive momentum-space quasi-PDFs, we will have to define a consistent subtraction scheme to remove all divergent terms as $\xi_z\to 0$, before the Fourier transformation.  Without this subtraction, the momentum-space quasi-PDFs are ill-defined at large $\tilde{x}$ region.

{\it Note added:} Two independent studies of the renormalization of quasi-PDFs appeared recently \cite{Ji:2017oey,Green:2017xeu}, in which the same conclusion is reached although  approaches are very different from ours.
